\section{Discussion}
\label{sec:discussion}
The results align with the design intent in one respect and correct it in another. Each optimization does produce its largest effect on the dimension it was built for. What the measurements do \emph{not} support is the expectation that the three, as a group, favour the lower-capability end of the hardware spectrum: shader pre-compilation pays on both discrete adapters and is not measurable on the integrated one, because on that part shader compilation is not what the first frame is spent on, while compressed textures combined with a size cap penalize only the integrated adapter, whose derived cube render target comes out of the same system memory the application is using. The useful generalization is not about hardware tiers: an optimization helps where the cost it removes is actually being paid, and establishing where that is requires measuring on the target rather than reasoning about its class.

For educational deployment the practically significant finding concerns loading rather than rendering. None of the three scenes was frame-limited in a way a student would notice, whereas the first rendered frame cost between 150 and 950 ms depending on configuration -- the interval a student experiences as the scenario failing to appear, and where the measurable headroom lies. For content a student opens, uses briefly and closes, removing half the wait before the first frame is worth more than improving a frame rate the display is already capping. The dirty-flag result separately illustrates the value of choosing the right metric: a phase-level speed-up alone would read as a win and a frame-level cost alone as a defect, whereas the two together describe what the mechanism actually is.

\subsection{Limitations of the current engine}
Two limitations are worth noting. The engine currently assumes Three.js as its backend, and the adapter layer is not yet generalized, which would be required for a WebGPU backend. Some features common in commercial engines -- including an entity-component-system architecture, GPU-driven culling, OffscreenCanvas rendering, geometry batching, and screen-relative level-of-detail (LOD) systems -- are not yet implemented.

